% Anexos del reporte técnico final, Reto AgroCebada FIRA 2026 (máximo 4 cuartillas).
% Compilar con: xelatex anexos.tex
\documentclass[10pt,letterpaper]{article}
\usepackage[margin=2.54cm]{geometry}
\usepackage{fontspec}
\setmainfont{TeX Gyre Termes}
\usepackage{polyglossia}
\setdefaultlanguage{spanish}
\usepackage{booktabs,array,tabularx,longtable}
\usepackage{fancyhdr}
\setlength{\headheight}{14pt}
\pagestyle{fancy}
\fancyhf{}
\fancyhead[L]{Los Bien Dotados}
\fancyhead[R]{Reto AgroCebada FIRA 2026}
\fancyfoot[C]{\thepage}
\setlength{\parindent}{0pt}
\setlength{\parskip}{3pt}
\setlength{\emergencystretch}{2em}
\usepackage{titlesec}
\titleformat{\section}{\normalsize\bfseries}{Anexo \Alph{section}.}{0.5em}{}
\titlespacing*{\section}{0pt}{8pt}{3pt}
\renewcommand{\arraystretch}{1.05}
\begin{document}
\begin{center}{\large\bfseries Anexos}\end{center}

\section{Actividades y aportaciones de cada integrante}
Las actividades marcadas con [por confirmar] son propuestas que el equipo debe validar antes de la entrega.

\small
\begin{tabularx}{\textwidth}{@{}>{\raggedright\arraybackslash}p{3.1cm}X@{}}
\toprule
Integrante & Actividades y aportaciones \\
\midrule
José Alán Pomares Valdés (modelado y análisis) &
Preparación y limpieza de los datos satelitales y climáticos e ingeniería de las 91 variables por parcela. Análisis de las etiquetas y hallazgo de los dos regímenes de registro. Diseño de la validación (esquemas A y B, validación anidada), de la regla de decisión y de las pruebas de permutación. Modelado con proceso gaussiano, rezago espacial e intervalos conformales. Capa de crédito y modelo municipal con el SIAP. Análisis de las preguntas guía. Redacción de los reportes técnicos. \\
\addlinespace
Vicente Josimar López Martínez (aplicación web) &
Desarrollo de la aplicación web con mapa de parcelas, predicción con intervalo, probabilidad bajo un umbral elegido por el usuario y ejecución del modelo en vivo. Levantamiento de necesidades de analistas y productores [por confirmar]. Pruebas de uso y corrección de errores de la interfaz [por confirmar]. Despliegue y publicación del enlace público [por confirmar]. Aportación de ideas en reuniones y revisión cruzada de resultados [por confirmar]. \\
\addlinespace
Luis Emmanuel Ramiro Luna (video) &
Elaboración del video de presentación del proyecto. Guion del video a partir de los resultados del reporte [por confirmar]. Grabación y edición [por confirmar]. Revisión de que el mensaje sea claro para productores y analistas [por confirmar]. Aportación de ideas en reuniones y revisión cruzada de resultados [por confirmar]. \\
\addlinespace
Luis Gustavo Olivares Soto (asesor) &
Revisión crítica de la metodología y de los reportes, con observaciones sobre la validación, el modelo municipal y la redacción. \\
\bottomrule
\end{tabularx}
\normalsize

\section{Uso de inteligencia artificial}
Usamos Claude Code, un asistente de modelos Claude de Anthropic, como apoyo de programación, auditoría de código y redacción. Lo empleamos para generar código base, auditar el modelo y los experimentos, preparar documentos y proponer redacciones. El equipo revisó cada resultado, lo contrastó con los datos y decidió qué incorporar. Las cifras del reporte salen de código ejecutado por el equipo. La tabla presenta síntesis formales de los prompts principales, no su texto original. Las fechas están en hora del centro de México.

\small
\begin{center}\textbf{Síntesis de los prompts principales}\end{center}
\vspace{-4pt}
\begin{longtable}{@{}r >{\raggedright\arraybackslash}p{1.2cm} >{\raggedright\arraybackslash}p{2.7cm} >{\raggedright\arraybackslash}p{4.1cm} >{\raggedright\arraybackslash}p{4.5cm}@{}}
\toprule
N & Fecha & Objetivo & Síntesis del prompt & Resultado y uso por el equipo \\
\midrule
\endhead
%%IA%%
\bottomrule
\end{longtable}
\normalsize

\par\vspace{8pt}\noindent\begin{minipage}{\textwidth}
\textbf{Anexo C. Predicciones de las 59 parcelas}\par\vspace{3pt}
Rendimiento predicho en t/ha con el intervalo del percentil 10 al 90 (p10 y p90).

\small
\setlength{\tabcolsep}{3pt}
\begin{tabular}{@{}llrrr@{\hspace{10pt}}llrrr@{}}
\toprule
ID & Municipio & Pred. & p10 & p90 & ID & Municipio & Pred. & p10 & p90 \\
\midrule
%%PRED%%
\bottomrule
\end{tabular}
\end{minipage}
\end{document}
